\section{Tercera iteración de la Fase 3: La marca de primer jugador en el diseñador}
\label{sec:PrimerJugador}

La iteración anterior (\emph{F3-I2}) cerró la respuesta al teclado, pero presentar el segmento inicial del \textbf{jugador activo} quedaba sujeto de una búsqueda: el motor tomaba el primer
segmento de tipo jugador que encontraba en la lista. No existía nada en el diseñador para indicar cual era el primer segmento a presentar al iniciar el juego \emph{por dónde empieza la partida}. Esta iteración retrocede
un paso al diseñador para dársela: un nuevo estado de los segmentos de tipo jugador, la marca de
\textbf{primer jugador}, con reglas sencillas y visibles en el panel de propiedades.

\subsection{El diseño}
\label{subsec:DisenoF3I3}

\begin{itemize}
	\item \textbf{Un checkbox en el panel}. Cuando el tipo del segmento es
	\emph{jugador}, el panel de propiedades muestra la fila «Jugador inicial»: marcarla
	dice que \emph{ese} es el segmento jugador con el que la partida empieza (el que el
	motor enciende al arrancar y del que salen los enlaces de dirección). Para los demás
	tipos la fila no existe: la marca es cosa del jugador, como sus enlaces.
	\item \textbf{Si no se marca ninguno}, el motor usa el primero que se encuentre en
	la lista de segmentos (la búsqueda de F3-I2 queda como salvavidas): el sistema
	funciona siempre, marcado o no.
	\item \textbf{Si se marcan dos, prevalece el que se acaba de marcar}. Solo puede
	haber una marca: marcar un segundo segmento le quita la marca al anterior. La
	intención del autor es siempre la última acción, y el sistema la respeta sin preguntar.
\end{itemize}


\begin{figure}[H]
	\centering
	\includegraphics[width=0.9\linewidth]{imagenes/propias/C05-I022}
	\caption[]{Marcar primer jugador.}
	\label{fig:c05-i022}
\end{figure}

\subsection{Cómo se realiza}
\label{subsec:RealizacionF3I3}

\begin{itemize}
	\item \textbf{Una propiedad más}. La marca vive como la propiedad
	\texttt{primero} del componente, junto a tipo, orden y enlaces: entra en el
	\texttt{propPorDefecto} (nada marcado por defecto) y en los \texttt{CAMPOS\_PROP} del
	historial, con lo que \textbf{viaja en el deshacer/rehacer por inversión} como
	cualquier otro campo, sobrevive a las uniones (el \emph{merge} de \texttt{reconstruir}
	la fusiona como el resto) y se reproduce al recargar la autoría.
	\item \textbf{Marcar dos es una acción atómica}. Quitarle la marca al anterior y
	dársela al nuevo se registran como el mismo bloque del historial (mismo número,
	\texttt{relacionados}): deshacer revierte \emph{ambas} cosas de un golpe y rehacer
	las repone. La unicidad no es un chequeo posterior: se mantiene en el momento del
	marcaje, recorriendo los grupos jugador marcados.
	\item \textbf{Cambiar de tipo descarta la marca}. Al dejar de ser jugador, la marca
	se descarta igual que los enlaces de dirección del tipo anterior (mismo aviso en
	consola, misma entrada del historial): nada queda huérfano en memoria ni en el GWI.
	\item \textbf{El GWI la dice solo si está}. El segmento marcado viaja con
	\texttt{primero: true} y los demás sin el campo (autoexplicativo y editable a mano;
	si un GWI editado a mano llevase dos marcas, el motor toma la primera que encuentra).
	La lista auxiliar de lo numerable muestra \textbf{«PRIMER JUGADOR: seg\_NNN»} junto a
	rutas y contadores: ver por dónde empieza la partida sin buscarlo en el mapa.
	\item \textbf{El motor solo cambia una prioridad}. \texttt{buscarJugador} consulta
	primero la marca del diseñador, y solo si no hay cae a la búsqueda anterior (tipo
	jugador, luego primer segmento con enlaces, luego el primero de la lista con aviso).
	La regla del usuario queda tal cual: \emph{la marca es la palabra del autor, la
		búsqueda es el salvavidas}.
\end{itemize}

\subsection{Verificación}
\label{subsec:VerificacionF3I3}

La simulación recorre el flujo completo ---diseñador, guardado/carga y motor--- con un
GWB de cuatro segmentos y un GWI de cuarenta:

\begin{itemize}
	\item \textbf{Diseñador}: clic en el mapa, tipo jugador y marca: el GWI en memoria
	lleva \texttt{primero: true}, la lista auxiliar lo muestra y la barra avisa; marcar un
	segundo le quita la marca al primero (prevalece el último); \textbf{deshacer} restaura
	la marca original y \textbf{rehacer} vuelve al segundo (acción atómica); desmarcar deja
	la lista sin marca con aviso; cambiar el tipo a juego descarta la marca con aviso en
	consola.
	\item \textbf{Round-trip}: un GWI cuya sección \texttt{autoria.cambios} contiene el
	marcaje se recarga y la reconstrucción reproduce la marca exactamente (la partida
	vuelve a empezar por donde el autor dijo).
	\item \textbf{Motor}: con la marca en el índice 5, el jugador activo es el 5
	(presentado a píxel; el segmento tipo jugador sin marca NO se presenta); sin marca,
	el motor cae al primero de la lista; los enlaces de dirección se siguen desde el
	segmento marcado igual que siempre. Cero errores de consola.
\end{itemize}

